\centerline{\bf \Huge Гомотетия посложнее}

\problem{1}
\textbf{Теорема Дезарга.} У треугольников $A_1 B_1 C_1$ и $A_2 B_2 C_2$ соответственные стороны попарно параллельны. Докажите, что прямые $A_1 A_2, B_1 B_2, C_1 C_2$ пересекаются в одной точке или параллельны.

\problem{2}
В треугольник $A B C$ вписана окружность $\omega$. Точки $M_a, M_b$ и $M_c$ - середины сторон $B C, C A$ и $A B$ соответственно. Точки $K_a, K_b$ и $K_c$ симметричны точкам касания окружности $\omega$ со сторонами $B C, C A$ и $A B$ относительно биссектрис углов $\angle A, \angle B$ и $\angle C$ соответственно. Докажите, что прямые $M_a K_a, M_b K_b$ и $M_c K_c$ пересекаются в одной точке.

\problem{3}
В треугольнике $A B C$ проведена биссектриса $A D$. Общая касательная $\ell$ к окружностям $(A B D)$ и $(A C D)$ касается их в точках $M$ и $N$ соответственно. Докажите, что прямая $\ell$ касается окружности, проходящей через середины отрезков $M N, B D, C D$.

\problem{4}
\textbf{Теорема о трех гомотетиях.} Докажите, что композиция гомотетий $H_{O_1}^{k_1}$ и $H_{O_2}^{k_2}$ является гомотетией $H_{O_3}^{k_3}$, где $k_3=k_1 k_2$ и точка $O_3$ лежит на прямой $O_1 O_2$.


\problem{5}
В треугольник $A B C$ вписана окружность $\omega$ с центром в точке $I$, касающаяся стороны $B C$ в точке $D$. Пусть $\omega_B$ и $\omega_C$ - вписанные окружности треугольников $A B D$ и $A C D$ соответственно, касающиеся стороны $B C$ в точках $E$ и $F$ соответственно. Пусть $P$ точка пересечения отрезка $A D$ с линией центров окружностей $\omega_B$ и $\omega_C, X$ - точка пересечения прямых $B I$ и $C P, Y$ - точка пересечения прямых $C I$ и $B P$. Докажите, что прямые $E X$ и $F Y$ пересекаются на окружности $\omega$.

\problem{6}
(a) \textbf{Лемма Архимеда.} Две окружности $\Omega$ и $\omega$ касаются в точке $A$. Прямая $\ell$ пересекает окружность $\Omega$ в точках $B$ и $C$ и касается окружности $\omega$ в точке $D$. Докажите, что прямая $A D$ проходит через середину дуги окружности $\Omega$, стягиваемую хордой $B C$ (разберите случаи внешнего и внутреннего касания). 

(b) Вокруг равнобедренного треугольника $A B C$ ( $A B=A C$ ) описана окружность $\Omega$. Окружность $\omega$ касается сторон $A B$ и $A C$ в точках $K$ и $L$ и касается внутренним образом окружности $\Omega$. Докажите, что центр вписанной окружности треугольника $A B C$ лежит на отрезке $K L$. 

(с) \textbf{Лемма Верьерра.} Докажите утверждение предыдущего пункта для произвольного треугольника.